\chapter{The Method of Steepest Descent}\label{Sec3.4}

The method of steepest descent, which is also commonly known as the saddle point method, is a technique for finding asymptotic approximations to integrals of the form
\begin{equation}\label{Eq:steepest_general}
    I(k) = \int_{C} f(z)\, e^{k\phi(z)}\, dz
\end{equation}
where $k$ is a large real parameter, $\phi(z)$ is an analytic function of the 
complex variable $z$, $f(z)$ is a further analytic function, and $C$ is some 
contour in the complex plane. The main idea is that as $k \to \infty$, the 
factor $e^{k\phi(z)}$ becomes exponentially large near the maximum of 
$\text{Re}(\phi(z))$ and exponentially small everywhere else. The entire integral is therefore controlled by a small neighbourhood of the point where $\text{Re}(\phi(z))$ is largest along the contour, and by deforming $C$ to pass through this point in a particular way the integral can be approximated to any desired order in $k$. \cite{RedBook,Asymptotic,Steepestdescent,Steepestorigin, bleistein1986,miller2006}

\subsection{Saddle Points}

The point $z_0$ is called a saddle point of $\phi$ if $\phi'(z_0) = 0$. Then more generally, suppose the first $n-1$ derivatives of $\phi$ all vanish at $z_0$, but the $n$-th derivative does not:$$\phi'(z_0) = \phi''(z_0) = \cdots = \phi^{(n-1)}(z_0) = 0, \qquad \phi^{(n)}(z_0) \neq 0$$Then the Taylor expansion of $\phi$ around $z_0$ begins at the $n$-th term, and the dominant term for correction is
\begin{equation}\label{Eq:taylor_saddle}    \phi(z) - \phi(z_0) \approx \frac{(z-z_0)^n}{n!}\,\phi^{(n)}(z_0)
\end{equation}
Writing $\phi^{(n)}(z_0)$ in polar form as $\phi^{(n)}(z_0) = |\phi^{(n)}(z_0)|e^{i\alpha}$, this then becomes$$\phi(z) - \phi(z_0) \approx \frac{(z-z_0)^n}{n!}\,|\phi^{(n)}(z_0)|\,e^{i\alpha}$$The name saddle point arises because since $\phi(z)$ is analytic, $\text{Re}(\phi(z))$ is harmonic and therefore cannot attain a true maximum or minimum in the interior of its domain. Instead the critical points of $\text{Re}(\phi)$ are saddle points, where it increases in some directions and 
decreases in others. \cite{RedBook, Asymptotic, Steepestdescent, Steepestorigin} 

\subsection{The Steepest Descent Directions}

Writing $z - z_0 = re^{i\theta}$ in polar form, where $r > 0$ is the distance from $z_0$ and $\theta$ is the direction of travel in the complex plane, the Taylor expansion \eqref{Eq:taylor_saddle} therefore becomes$$\phi(z) - \phi(z_0) \approx \frac{r^n}{n!}\,|\phi^{(n)}(z_0)|\,e^{i(\alpha + n\theta)}$$For $e^{k\phi(z)}$ to decay away from $z_0$ we require $\text{Re}(\phi(z) - \phi(z_0)) < 0$, that is, the above expression must be real and negative. Since $\frac{r^n}{n!}|\phi^{(n)}(z_0)|$ is always a positive real number, this condition requires $e^{i(\alpha + n\theta)} = -1$, which holds precisely when$$\alpha + n\theta = (2m+1)\pi, \qquad m = 0, 1, \ldots, n-1$$Solving for $\theta$ gives exactly $n$ steepest descent directions:
\begin{equation}\label{Eq:descent_directions}
    \theta = \frac{(2m+1)\pi - \alpha}{n}, \qquad m = 0, 1, \ldots, n-1
\end{equation}
These are the directions in which moving away from $z_0$ causes the integrand to decay most rapidly. \cite{RedBook,Asymptotic,Steepestdescent,Steepestorigin, olver1997}

\subsection{The Change of Variables to Laplace Form}

Once the contour $C$ has been deformed, using Cauchy's theorem, to pass through $z_0$ along one of the steepest descent directions $\theta$, we introduce the real substitution
\begin{equation}\label{Eq:substitution}
    \phi(z) - \phi(z_0) = -t, \qquad t \geq 0
\end{equation}where $t$ starts at $0$ when $z = z_0$ and increases as we move away from $z_0$ along the deformed contour. The key point is that by construction of the steepest descent path, $\phi(z) - \phi(z_0)$ is real and negative along the contour, so this substitution is well defined and $t$ is a genuine non-negative real variable.
\\
Substituting \eqref{Eq:substitution} into $e^{k\phi(z)}$ gives$$e^{k\phi(z)} = e^{k[\phi(z_0) + (\phi(z) - \phi(z_0))]} = e^{k\phi(z_0)}\,e^{-kt}$$so the complicated complex exponential is replaced by a simple real decaying exponential $e^{-kt}$. Changing the integration variable from $z$ to $t$ via $dz = \frac{dz}{dt}\,dt$, the integral becomes
\begin{equation}\label{Eq:laplace_form}
    I(k) = e^{k\phi(z_0)}\int_0^\infty f(z(t))\,e^{-kt}\,\frac{dz}{dt}\,dt
\end{equation}
which is a standard Laplace-type integral in the variable $t$, amenable to asymptotic expansion for large $k$ via Watson's lemma. \cite{olver1997,Steepestdescent}
\\
To compute $\frac{dz}{dt}$ explicitly, we use the Taylor expansion \eqref{Eq:taylor_saddle} and the substitution \eqref{Eq:substitution} together:$$(z-z_0)^n \approx \frac{-n!\,t}{|\phi^{(n)}(z_0)|\,e^{i\alpha}}$$Taking the $n$-th root and selecting the branch corresponding to the descent direction $\theta$ gives $$z - z_0 \approx \left(\frac{n!\,t}{|\phi^{(n)}(z_0)|}\right)^{1/n} e^{i\theta}$$ Differentiating with respect to $t$:
\begin{equation}\label{Eq:jacobian}
    \frac{dz}{dt} \approx \frac{1}{n}\left(\frac{n!}{|\phi^{(n)}(z_0)|}
    \right)^{1/n} t^{1/n - 1}\,e^{i\theta}
\end{equation}

\subsection{Extending the Upper Limit to Infinity}

The integral \eqref{Eq:laplace_form} has been written with upper limit $\infty$, but in practice the contour may only extend to some finite point, so the upper limit in the $t$ variable is some finite value $T$ rather than $\infty$. The justification for replacing $T$ by $\infty$ is that the tail of the integral beyond $T$ is exponentially negligible compared to the leading term. to be precise, since $\phi(z) - \phi(z_0) = -t < 0$ along the descent path and $t$ grows as we move away from $z_0$, the factor $e^{-kt}$ decays exponentially for $t > T$. For large $k$ the contribution from $t \geq T$ is of order $e^{-kT}$, which is smaller than any power of $1/k$ and is therefore asymptotically negligible. The upper limit may therefore be extended to $\infty$ without affecting the asymptotic expansion to any order in $k$, and the estimation lemma \ref{Eq3.1.7} provides the rigorous justification for this step.

\subsection{Interior and Boundary Saddle Points}

The framework above assumes the saddle point $z_0$ lies in the interior of the deformed contour, so that the steepest descent path passes through $z_0$ from both sides and the Laplace-type integral \eqref{Eq:laplace_form} runs from $-\infty$ to $+\infty$ through $z_0$. In this case the leading order contribution is a full Gaussian integral$$\int_{-\infty}^{\infty} e^{-kt}\,dt$$and the asymptotic formula \eqref{Eq:steepest_result} applies directly.

However, if $z_0$ is a \textit{boundary saddle point}, meaning it lies at an endpoint of the contour rather than in its interior, then the steepest descent path leaves $z_0$ in only one direction and the integral \eqref{Eq:laplace_form} runs from $0$ to $+\infty$ rather than from $-\infty$ to $+\infty$. In this case only half of the Gaussian integral contributes, and the leading order result is reduced by a factor of $\frac{1}{2}$. Concretely, for the standard case $n = 2$, $\beta = 1$ the formula \eqref{Eq:steepest_standard} becomes
\begin{equation}\label{Eq:steepest_boundary}
    I(k) \sim \frac{1}{2}\,f(z_0)\,e^{k\phi(z_0)}\,e^{i\theta}
    \sqrt{\frac{2\pi}{k\,|\phi''(z_0)|}}
    = f(z_0)\,e^{k\phi(z_0)}\,e^{i\theta}
    \sqrt{\frac{\pi}{2k\,|\phi''(z_0)|}}
    \qquad \text{as } k \to \infty
\end{equation}
It is therefore important when applying the method to identify whether the saddle point is interior or at a boundary, as the two cases give different prefactors. \cite{olver1997,bleistein1986}

\subsection{The General Asymptotic Formula}

We assume $f(z)$ behaves near $z_0$ as a power law of the form$$f(z) \sim f_0\,(z - z_0)^{\beta - 1}$$for some constants $f_0$ and $\beta$. This accommodates the case where $f$ has a branch point or zero at $z_0$. When $f$ is an ordinary analytic function with $\beta = 1$ this reduces to $f(z) \approx f_0 = f(z_0)$. Substituting the expression for $z - z_0$ gives $$f(z) \sim f_0\left(\frac{n!\,t}{|\phi^{(n)}(z_0)|}\right)^{(\beta-1)/n} e^{i\theta(\beta-1)}$$Collecting $f(z)$, $e^{-kt}$, and $\frac{dz}{dt}$ inside the integral, the powers of $t$ combine as $$\frac{\beta - 1}{n} + \frac{1}{n} - 1 = \frac{\beta}{n} - 1$$ and the phase factors combine as $e^{i\theta(\beta-1)}\cdot e^{i\theta} = e^{i\beta\theta}$. The integral \eqref{Eq:laplace_form} therefore becomes $$I(k) \sim e^{k\phi(z_0)}\cdot\frac{f_0\,e^{i\beta\theta}}{n}\left(\frac{n!}{|\phi^{(n)}(z_0)|}\right)^{\beta/n}\int_0^\infty t^{\beta/n - 1}\,e^{-kt}\,dt$$
Evaluating the remaining integral using the Gamma function via the substitution $u = kt$:$$\int_0^\infty t^{\beta/n - 1}\,e^{-kt}\,dt = k^{-\beta/n}\int_0^\infty u^{\beta/n - 1}\,e^{-u}\,du = k^{-\beta/n}\,\Gamma\!\left(\frac{\beta}{n}\right)$$and so the general asymptotic formula for an interior saddle point is
\begin{equation}\label{Eq:steepest_result}
    I(k) \sim \frac{f_0\,(n!)^{\beta/n}\,e^{i\beta\theta}\,e^{k\phi(z_0)}\,
    \Gamma\!\left(\frac{\beta}{n}\right)}{n\left(k\,|\phi^{(n)}(z_0)|\right)^{\beta/n}}
    \qquad \text{as } k \to \infty
\end{equation}
For a boundary saddle point the right hand side of \eqref{Eq:steepest_result} 
is multiplied by $\frac{1}{2}$. If there are multiple saddle points through 
which the deformed contour passes, each contributes a term of this form and the 
contributions are summed to give the full asymptotic expansion. \cite{RedBook,Asymptotic, Steepestdescent,Steepestorigin}

\subsection{The Standard Case}

The most commonly encountered situation is $n = 2$ and $\beta = 1$, where $z_0$ 
is a simple saddle point and $f(z)$ is analytic and non-zero at $z_0$. In this 
case $\Gamma(1/2) = \sqrt{\pi}$ and \eqref{Eq:steepest_result} reduces to
\begin{equation}\label{Eq:steepest_standard}
    I(k) \sim f(z_0)\,e^{k\phi(z_0)}\,e^{i\theta}
    \sqrt{\frac{2\pi}{k\,|\phi''(z_0)|}}
    \qquad \text{as } k \to \infty
\end{equation}
where $\theta$ is the angle of the steepest descent path through $z_0$, 
determined by $\cos(2\theta + \alpha) < 0$ with $\alpha = \arg(\phi''(z_0))$.

\subsection{Summary of the Method}

In practice the method proceeds as follows. First, identify the saddle points 
of $\phi(z)$ by solving $\phi'(z) = 0$ and determine the order $n$ of each. 
Second, determine which saddle points the contour $C$ can be deformed through 
without crossing any singularities of the integrand, using Cauchy's theorem to 
justify the deformation, and identify whether each relevant saddle point is 
interior or at a boundary of the contour. Third, at each relevant saddle point 
find the steepest descent directions $\theta$ from \eqref{Eq:descent_directions}. 
Fourth, apply the substitution $\phi(z) - \phi(z_0) = -t$ to convert the 
integral to Laplace form \eqref{Eq:laplace_form}, verify that the tail beyond 
any finite endpoint is exponentially negligible, and extend the upper limit to 
$\infty$. Finally, apply \eqref{Eq:steepest_result} at each saddle point, 
including the factor of $\frac{1}{2}$ for any boundary saddle point, and sum 
the contributions to obtain the leading order asymptotic behaviour of $I(k)$ 
as $k \to \infty$.



\subsection{Problems}
\textbf{Question 1:} From the 'Complex Variables' book \cite{RedBook}, Find the asymptotic behaviour as $k\to\infty$ of the integral 
\begin{equation}
    I(k)=\int^{\infty}_{-\infty}e^{k(iz-z^2)}
\end{equation}
First step is to identify the $\phi(z)$ and $f(z)$, therefore looking at the standard form $I(k)=\int_C f(z)e^{k\phi(z)}dz$ we can then see that $f(z)=1$ and $\phi(z)=iz-z^2$ 
\\
Now we can work out the saddlepoints, saddlepoints are where $\phi'(z)=0$ therefore the differential of $\phi(z)$ is $\phi'(z)=i-2z$, so by setting this to 0 we will then find the saddlepoint. $$i-2z=0 \Rightarrow z_0=\frac{i}{2}$$ So there is a saddlepoint at $z_0=\frac{i}{2}$ which lies on the imaginary axis. Now we can check the second derivate to determine whether it is a simple pole (n=2). $\phi''(z)=-2\neq0$. Since this in not equal to zero this is a simple saddlepoint of order $n=2$. 
\\
Next is to evaluate the saddlepoint as $\phi(z_0)$ is needed for the final formula. Substitute $z_0=\frac{i}{2}$ $$\phi(z_0)=i\cdot\left(\frac{i}{2}\right)-\left(\frac{i}{2}\right)^2=-\frac{1}{4}$$ So $e^{k\phi(z_0)}=e^{-\frac{k}{4}}$ this is the exponential factor which will be in final result. 
\\
Next what is needed is to find the steepest descent direction, therefore what is needed is to find the directions which to travel away from $z_0$ so that the integral decays. We can then write 
\begin{equation}
    \phi''(z_0)=-2=2e^{i\pi}
\end{equation}
So 
\begin{equation}
    |\phi''(z_0)|=2 \text{ and } \alpha=\arg(\phi''(z_0))=\pi
\end{equation}
Then the steepest descent directions for a simple saddlepoint $n=2$ come from the condition 
\begin{equation}
    \alpha+2\theta=(2m+1)\pi \ \ \ \ m=0,1
\end{equation}
Now we can substitute $\alpha=\pi$ 
\begin{equation}
    \pi+2\theta=(2m+1)\pi \Rightarrow 2\theta=2m\pi\Rightarrow \theta=m\pi
\end{equation}


So the two directions of descent are $\theta=0$ which points right along the real axis and $\theta=\pi$ which points left along the real axis, together these form a horizontal path through $z_0$, which makes sense as the steepest descent path here is the horizontal line $\Im(z)=\frac{1}{2}$. 
\\
The next step is to deform the contour, as the original contour is the real axis but we need it to pass through $z_0=\frac{i}{2}$ along the steepest descent path, which is the horizontal line 
\begin{equation}
    z=x+\frac{i}{2} \ \ \ \ \text{where } x\in(-\infty,\infty)
\end{equation}
\\

\begin{figure}[h]
\centering
\begin{tikzpicture}[
    scale=1.0,
    >=Stealth,
    arr/.style={
        decoration={markings, mark=at position 0.5 with 
        {\arrow[line width=1.2pt]{>}}},
        postaction=decorate
    },
]

% Axes
\draw[->, line width=0.7pt] (-2.0,0) -- (2.0,0) 
    node[right] {$\mathrm{Re}(z)$};
\draw[->, line width=0.7pt] (0,-0.4) -- (0,2.0) 
    node[above] {$\mathrm{Im}(z)$};

% Original contour: real axis
\draw[line width=1.5pt, arr, dashed] (-1.8,0) -- (1.8,0);
\node[below, font=\small] at (-1.2,-0.12) {$C$};

% Deformed contour: horizontal line at Im(z) = 1
\draw[line width=1.5pt, arr] (-1.8,1.0) -- (1.8,1.0);
\node[below, font=\small] at (-1.2,1.05) {$C'$};

% Saddle point
\filldraw (0,1.0) circle (2pt);
\node[above left, font=\small] at (0,1.0) {$z_0 = \frac{i}{2}$};

% Deformation arrow
\draw[->, line width=0.8pt] (1.4,0.08) -- (1.4,0.92);
\node[right, font=\small] at (1.45,0.5) {deform};

% Origin
\node[below left, font=\small] at (0,0) {$O$};

\end{tikzpicture}
\caption{Contour deformation for $I(k)$. The original contour $C$ along 
the real axis is deformed to $C'$ passing through the saddle point $z_0 = i/2$.}
\label{fig:steepest_contour}
\end{figure}

\\
Then by Cauchy's theorem, since $f(z)e^{k\phi(z)}$ has no singularities we can freely deform the contour from the real axis to the horizontal line without changing the value of the integral 
\\
Now we can apply the substitution $\phi(z)-\phi(z_0)=-t$. Then verify this works with the deformed contour. Write $z=x+\frac{i}{2}$ so $z-z_0=x$. Then 
\begin{equation}
    \phi(z)-\phi(z_0)=\phi\left(x+\frac{i}{2}\right)-\phi\left(\frac{i}{2}\right)
\end{equation}
$$\phi\left(x+\frac{i}{2}\right)=i\left(x+\frac{i}{2}\right)-\left(x+\frac{i}{2}\right)^2=-\frac{1}{4}-x^2$$ Therefore $$\phi(z)-\phi(z_0)=-\frac{1}{4}-x^2-\left(-\frac{1}{4}\right)=-x^2$$ Therefore 
\begin{equation}
    \phi(z)-\phi(z_0)=-x^2
\end{equation}
Which is real and negative for all $x\neq0$ this confirms the contour is a valid steepest descent path setting $t=x^2$, so we have 
\begin{equation}
    e^{k\phi(z)}=e^{k\phi(z_0)}e^{-kx^2}
\end{equation} 
\\
Now we can finally transform the integral, on the deformed contour $z=x+\frac{i}{2}$ we have $dz=dx$ the integral becomes 
\begin{equation}
    I(k)=\int^{\infty}_{-\infty}e^{k\phi(z_0)}e^{-kx^2}dx=e^{-\frac{k}{4}}\int^{\infty}_{\infty}e^{-kx^2}dx
\end{equation}
This is now in the form of a standard Gaussian integral. The saddlepoint is the interior, so there is no factor $\frac{1}{2}$. Now after evaluating this Gaussian integral using the standard result $\int^{\infty}_{-\infty}e^{-kx^2}dx=\sqrt{\frac{\pi}{k}}$ 
\begin{equation}
    I(k)\sim e^{-\frac{k}{4}}\sqrt{\frac{\pi}{k}}
\end{equation}
\\
\\
\textbf{Problem 2:} From the 'Complex Variables' book \cite{RedBook}, Use the method of steepest descent to find the leading asymptotic behaviour as $k\to\infty$ of 
\begin{equation}
    \int^{\infty}_{-\infty}\frac{te^{ik(\frac{t^3}{3}+t)}}{1+t^4}dt
\end{equation}
\\
Again $I(k)=\int_Cf(z)e^{k\phi(z)}dz$ so it can be seen that $f(
z)=\frac{z}{1+z^4}$ and $\phi(z)=i\left(\frac{z^3}{3}+z\right)$ therefore as the saddlepoints occur when $\phi'(z)=0$ $$\phi'(z)=(z^2+1)i=0\Rightarrow z^2=-1\Rightarrow z=\pm i$$ Both of these are simple saddlepoints since $\phi''(z)=2iz\neq0$ which confirms the order is $n=2$ 
\\
Now to evaluate the saddlepoints at both $z=\pm i$: 
\\
At $z_0=i$\\
$\phi(i)=i\left(\frac{i^3}{3}+i\right)=i\left(\frac{2i}{3}\right)=\frac{-2}{3}$
\\
At $z_0=-i$\\
$\phi(-i)=i\left(\frac{(-i)^3}{3}-i\right)=i\left(\frac{-2i}{3}\right)=\frac{2}{3}$
\\
Since $e^{k\phi(-i)}=e^{\frac{2k}{3}}$ grows exponentially as $k\to\infty$ the saddlepoint at $z=-i$ cannot lie on a valid steepest descent contour for any large positive $k$. Therefore only $z=+i$ is relevant. 
\\
$\phi''(i)=2i\cdot i=-2$ so $|\phi''(z_0)|=2$ and $\alpha=\arg(-2)=\pi$ therefore using \ref{Eq:descent_directions} using $$\alpha+2\theta=(2m+1)\pi\Rightarrow\pi+2\theta=(2m+1)\pi$$ $$\Rightarrow 2\theta=2m\pi\Rightarrow \theta=m\pi$$ where $m=0,1$ so $\theta=0,\pi$, the steepest descent path through $z_0=i$ is the horizontal line $\Im(z)=1$ so $z=x+i$ for $x\in(-\infty,\infty).$ 
\\
We then verify the deformation of the contour is valid. The poles of $f(z)=\frac{z}{1+z^4}$ are where $z^4=-1$ giving 
\begin{equation}
    z=e^{i\pi(2m+1)/4} \ \ \ \ \text{ where } m=0,1,2,3
\end{equation}
\begin{equation}
    \Rightarrow z=\frac{\pm1\pm i}{\sqrt{2}}
\end{equation}
All four poles have $|\Im(z)|=\frac{1}{\sqrt{2}}<1$ so none of them lie between the the real axis and $\Im(z)=1$> There by Cauchy's theorem the contour can be shifted upwards without crossing a singularity so the deformation is valid. 
\\
Now verifying the steepest descent path: write $z=x+i$ so $z-z_0=x$ 
\begin{equation}
    \phi(z)-\phi(z_0)=i\left(\frac{(x+i)^3}{3}+(x+i)\right)-\left(\frac{-2}{3}\right)
\end{equation}
$(x+i)^3=x^3+3x^2i-3x-i$ $$\phi(z)-\phi(z_0)=i\left(\frac{x^3-3x}{3}+\frac{3x^2-i}{3}+x+i\right)+\frac{2}{3}$$ $$=\frac{ix^3}{3}-x^2-\frac{2}{3}+\frac{2}{3}=\frac{ix^3}{3}-x^2$$ For the steepest descent approximation, we need this to be real and negative near $z_0$. The $-x^2$ term will dominate when $x$ is small and is indeed real and negative. The imaginary term will only contribute as an oscillatory phase. therefore the path is a valid descent path. 
\\
Applying the standard case formula \ref{Eq:steepest_standard} with $n=2, \beta=1, f(z)$ analytic and nonzero at $z_0=i$, $f(z_0)=\frac{i}{2}$ using
\begin{equation}
    I(k)\sim f(z_0)e^{k\phi(z_0)}e^{i\theta}\sqrt{\frac{2\pi}{k|\phi''(z_0)|}} \ \ \ \ k\to\infty
\end{equation} 
its shown that 
\begin{equation}
    I(k)\sim \frac{i}{2}\cdot e^{\frac{-2k}{3}}\cdot1\sqrt{\frac{2\pi}{2k}}
\end{equation}
\begin{equation}
    I(k)\sim \frac{i\sqrt{\pi}}{2\sqrt{k}}e^{\frac{-2k}{3}} \text{ as } k\to\infty
\end{equation}
\\
\textbf{Problem 3:} From the 'Complex Variables' book \cite{RedBook} Show that for $0\leq\theta\leq\frac{\pi}{2}$ that $$\int^{\theta}_{0} e^{-k\sec x} dx \sim \sqrt{\frac{\pi}{2k}}e^{-k}$$ 
As usual the first step is write it in standard for to evaluate, as $I(k)=\int_C f(z)e^{k\phi(z)}dz$ we can see again that $f(z)=1$ and by identifying $$e^{-k\sec x}=e^{k(-\sec x)}$$ so $\phi(z)=-\sec z$, and the original contour goes from $\theta=0$ to $\theta=\frac{\pi}{2}$. 
Again a saddlepoint is $z_0$ where $\phi'(z_0)=0$. So we then differentiate $\phi(z)=-\sec z$, $$\phi'(z)=-\frac{d}{dz}\sec z=-\sec z \tan z$$ then by setting $\phi'(z)=0$ $$-\sec z \tan z=0$$ Since $\sec z=\frac{1}{\cos z}$ is never equal to zero we need $\tan z=0$ therefore $$\tan z=0 \ \text{ when } z=0,\pm\pi,\pm2\pi,...$$. 
The only saddlepoint here is when $z_0=0$ as the rest are outside of our contours interval $[0,\frac{\pi}{2}]$, also note that $z_0=0$ will lie on the left side endpoint of the integration interval. 
Next is to classify the saddlepoint, so at $z_0=0$ we compute $\phi(z_0)$ and $\phi''(z_0)$. $$\phi(0)=-\sec(0)=-\frac{1}{\cos(0)}=-1$$ Now the second derivative of $\phi(z)$ $$\phi''(z)=\frac{d}{dz}\phi(z)=\frac{d}{dz}(-\sec z \tan z)=-\sec z \tan^2z-\sec^3z$$
At $z_0=0$ $\sec(0)=1$ and $\tan(0)=0$ so $$\phi''(0)=-(1)(0)^2-(1)^3=-1$$ It can then be seen that $\phi''(z_0)\neq0=-1$ which means the order of the saddlepoint is equal to $n=2$ and it is also real and negative which is what we need. 
Now near the saddlepoint $z_0=0$ we can write $z=re^{i\theta}$ for small $r>0$ and direction angle $\theta$, the taylor expansion gives, 
\begin{align*}
\phi(z) &\approx\phi(0)+\frac{1}{2}\phi''(0)(z-0)^2=-1+\frac{1}{2}(-1)r^2e^{2i\theta} \\
\\
\phi(z) &\approx-1-\frac{r^2}{2}e^{2i\theta}
\end{align*}
for the steepest descent path we need $f(z)-f(0)$ to be real and negative so that $e^{k(\phi(z)-\phi(0))}$ decays away from $z_0$ which requires $\frac{r^2}{2}e^{2i\theta}$ to be real and negative. Since $-\frac{r^2}{2}<0$ we need $e^{2i\theta}$ to be real and positive.
$$2\theta=0, \text{ or } 2\theta=2\pi \Rightarrow \theta=0,\pi$$ 
The next part would normally be where we deform the contour however since our original contour runs along the positive real axis from $0\to\frac{\pi}{2}$, the steepest descent direction $\theta=0$ is already aligned with the original contour. This means there is no contour deformation needed as the original path of integration is the steepest descent path. 
Along $\theta=0$ (the positive real axis) near $z_0=0$ $$\phi(z)\approx-1-\frac{r^2}{2}$$ So $e^{k\phi(z)}\approx e^{-k}\cdot e^{\frac{-kr^2}{2}}$, the factor $e^{\frac{-kr^2}{2}}$ decays like a Gaussian away from $r=0$ for large $k$. This confirms that for large $k$ the integrand is dominated by the region near $z=0$, and contributes from $x$ away from $0$ are exponentially small by comparison. 
\\
Since $z_0=0$ lies at the left endpoint of the contour rather than at an interior point, we only integrate over half of the steepest descent path, so from $z_0$ outward in the direction of $\theta=0$ rather than through $z_0$ in both directions. This means the final formula will acquire a factor of $\frac{1}{2}$.
\\ 
Now applying the standard case formula we get $$I(k)\sim \frac{1}{2}\cdot f(z_0)\cdot e^{k\phi(z_0)}\cdot e^{i\theta}\sqrt{\frac{2\pi}{k|\phi''(z_0)|}}$$
as $f(z_0)=1$, $e^{k\phi(z_0)}=e^{-k}$ and $e^{i\theta}=1$ the formula then simplifies to 
$$I(k)\sim\frac{1}{2}e^{-k}\sqrt{\frac{2\pi}{k}}\Rightarrow\sqrt{\frac{\pi}{2k}}e^{-k}$$